\section*{Limitations}
\label{sec:limitations}
% ACL requires an unnumbered "Limitations" section; it does NOT count toward the
% page limit. Be candid: reviewers reward honesty here.

\paragraph{Scale.} The human set has 29 labeled queries and only 13 error steps.
Bootstrap CIs for \texttt{caught\_error} are therefore uninformative (Claude's
reaches the boundary $1$); we report McNemar/Holm $p$-values and point estimates
as directional, not precise effect sizes. The 13 errors are also clustered
rather than spread evenly: they fall in only 5 of the 29 queries, with 6 in Q09
and 4 in Q22 (Appendix~\ref{app:human-errors}). Because step-level McNemar
treats each step as an independent trial, it overstates the effective sample
size for \texttt{caught\_error}; the 13 steps are closer to 5 (correlated)
observations. As a coarser, query-level check, a query counts as ``caught'' if
a grader flags \emph{any} of its error steps: \texttt{claude} catches at least
one error in $5/5$ error-containing queries, \texttt{gpt-oss} in $2/5$, and the
symbolic verifier in $0/5$, the same ranking as the step-level results. We
therefore treat the step-level $p$-values as directional evidence corroborated
by, rather than independent of, this coarser query-level comparison.

\paragraph{Judge coverage.} We test two judges from two labs. This is stronger
than a single judge, but it is not ``all judges''; we therefore claim
model-independence of the \emph{qualitative} pattern, not universal
generalization.

\paragraph{Annotation.} The ground truth is single-annotator: no second coder
labeled any subset for this submission, so we report no inter-coder agreement
statistic \citep{artstein2008agreement}. We mitigate by (i)~publishing the
labeling rubric, including the notation-strictness rule (a step is marked
incorrect if boundary/IBP notation is written wrong even when the final result is
correct), and (ii)~releasing all raw per-step labels so others can audit or
re-annotate. Absolute agreement rates should therefore be read cautiously; the
comparative pattern across graders is the primary claim. Separately, LLM judges
were graded against this rubric without being shown it (\S\ref{sec:protocol}), so
human-set precision/recall are agreement scores, not oracle-matched accuracy. The
annotator is also the author of the grading pipeline and the labeling rubric; we
mitigate by releasing the rubric and all raw labels, and by reporting the
in-fragment/out-of-fragment classification in Appendix~\ref{app:human-errors} so
readers can assess whether the rubric systematically disadvantages symbolic
grading.

\paragraph{Synthetic errors.} Injection mutations are programmatic and may not
match the real distribution of LLM mistakes; they provide controlled recall/FPR
measurements, not a naturalistic error sample. Their 9 base derivations are also
hand-authored textbook results, so they are cleaner and shorter than the
assistant's own traces, which is what makes the mutation the sole source of
error but also limits how far injection numbers transfer to real output.

\paragraph{Single domain / frozen generator.} All human-labeled derivations come
from one signals-and-systems corpus and one generator,
\texttt{qwen-3-235b-a22b-instruct-2507} (Cerebras), which the provider has since
deprecated. This is not a domain-transfer or generator-robustness study. A
retrieval ablation (or any regeneration) would require swapping in a live
replacement model and would break comparability with the frozen main study; we
leave that to future work.
